\documentclass[10pt, a4paper]{article}

% Packages
\usepackage[utf8]{inputenc}
\usepackage[T1]{fontenc}
\usepackage[german]{babel}
\usepackage{xcolor}
\usepackage{geometry}
\usepackage{hyperref}
\usepackage{charter}

% Layout settings
\geometry{left=2.5cm, right=2.5cm, top=2cm, bottom=2cm}
\pagestyle{empty}

% Colors
\definecolor{primarycolor}{RGB}{0, 51, 102}
\hypersetup{colorlinks=true, urlcolor=primarycolor, linkcolor=primarycolor}

\begin{document}

% SENDER INFO
\noindent
\textbf{\large \color{primarycolor} Kartavya Niraj Dikshit} \\
Halbmondstraße 7 \\
91054 Erlangen \\
E-Mail: \href{mailto:kartavya.dikshit@gmail.com}{kartavya.dikshit@gmail.com} \\
Tel: +49 15510940829 \\
LinkedIn: \href{http://www.linkedin.com/in/kartavya-dikshit}{linkedin.com/in/kartavya-dikshit} \\
GitHub: \href{https://github.com/KartavyaDikshit}{github.com/KartavyaDikshit}

\vspace{1.5cm}

% RECIPIENT INFO
\noindent
Siemens AG \\
Foundational Technologies (FT) \\
Recruiting Team \\
Lichtenbergstraße 1 \\
85748 Garching bei München

\hfill Erlangen, \today

\vspace{1cm}

% SUBJECT
\noindent
\textbf{\large \color{primarycolor} Bewerbung als Werkstudent (f/m/d) Leveraging Generative AI along the Software Development Lifecycle (SDLC)} \\
\textbf{Job-ID: 496890}

\vspace{0.8cm}

% SALUTATION
\noindent
Sehr geehrte Damen und Herren,

\vspace{0.4cm}

% CONTENT
\noindent
mit großer Begeisterung verfolge ich die Pionierarbeit von Siemens im Bereich der industriellen Digitalisierung. Als Masterstudent der Data Science an der Friedrich-Alexander-Universität Erlangen-Nürnberg (FAU) mit einer Spezialisierung auf Künstliche Intelligenz und Machine Learning suche ich eine Herausforderung, bei der ich meine Expertise in Generativer KI direkt in die Optimierung von Softwareentwicklungsprozessen einbringen kann. Die ausgeschriebene Stelle bei Siemens Foundational Technologies bietet hierfür die ideale Schnittstelle zwischen Forschung und Anwendung.

\vspace{0.3cm}
\noindent
In meinem bisherigen Werdegang habe ich mich intensiv mit der Architektur und Skalierung von KI-Modellen auseinandergesetzt. Während meiner Zeit bei Hidden Brains konnte ich die Latenzzeit von KI-gestützten Business-Intelligence-Systemen um 52\,\% senken. Besonders relevant für Ihre Anforderungen im SDLC ist meine Erfahrung in der Entwicklung von Multi-Agenten-Systemen (z.B. mit LangGraph) für automatisierte Code-Reviews sowie der Implementierung von RAG-basierten Dokumentationsassistenten. Diese Projekte haben mir gezeigt, wie LLMs nicht nur Code generieren, sondern die gesamte Qualitätssicherung und Wissensverteilung in einem Team transformieren können.

\vspace{0.3cm}
\noindent
Mein Studium an der FAU vertieft mein Verständnis für Deep-Learning-Architekturen und verteiltes Rechnen – Kompetenzen, die für die bei Siemens angestrebte Effizienzsteigerung in der Softwaretechnik essenziell sind. Ich bin sicher im Umgang mit Python, PyTorch und TensorFlow und verstehe es, komplexe Problemstellungen analytisch zu durchdringen. Meine Erfahrung bei Icertis hat zudem mein Verständnis für professionelle SDLC-Prozesse, Debugging und Code-Qualität in agilen Umgebungen geschärft.

\vspace{0.3cm}
\noindent
Siemens FT steht für Innovation auf höchstem Niveau. Ich bin hochmotiviert, mein Wissen über Prompt Engineering und Fine-Tuning einzusetzen, um gemeinsam mit Ihren Experten die nächste Generation KI-gestützter Entwicklungstools zu gestalten. Ich verfüge über sehr gute Englischkenntnisse (C1) und arbeite kontinuierlich an der Perfektionierung meines Deutschen (aktuell B1), um mich nahtlos in Ihr hybrides Team in Garching einzufügen.

\vspace{0.3cm}
\noindent
Gerne möchte ich Sie in einem persönlichen Gespräch von meiner fachlichen Eignung und meiner Begeisterung für diese Aufgabe überzeugen. Ich stehe Ihnen für 20 Stunden pro Woche zur Verfügung.

\vspace{0.8cm}

% CLOSING
\noindent
Mit freundlichen Grüßen

\vspace{0.8cm}
\noindent
Kartavya Niraj Dikshit

\end{document}